% Pregunta 4 (2025-marzo-A, n.º 4): descomposición de Hicks de una subida de p_x con un bien normal.
% U = (xy)^(1/2), m = 4, p_y = 1; p_x pasa de 1 a 2 (a escala, eje x con 1,3 cm por unidad).
%   A = (2, 2) en U0: xy = 4; C = (1, 2) en U1: xy = 2.
%   Recta compensada y = 4·sqrt(2) - 2x, tangente a U0 en B = (sqrt 2, 2·sqrt 2).
%   ES = x_B - x_A < 0 y ER = x_C - x_B < 0: los dos reducen x, luego x2 < x1.
% Autor: Víctor Gutiérrez Marcos
\documentclass[tikz,border=6pt]{standalone}
\input{tikz-tcee}
\begin{document}
\begin{tikzpicture}[tcee,x=1.3cm,y=1cm]
  \def\xA{2}\def\yA{2}\def\xB{1.41421}\def\yB{2.82843}\def\xC{1}\def\yC{2}
  \node[font=\normalsize\bfseries,text=tceemorado] at (2.3,6.1) {Bien normal: $x_2<x_1$};
  \ejes{$x$}{$y$}{4.7}{5.8}
  % Rectas de balance
  \draw[curva1] (0,4) -- (4,0) node[etiqueta,above right,text=tceemorado] {$p_x^0$};
  \draw[curvanueva] (0,4) -- (2,0) node[etiqueta,above right,text=tceerojo] {$p_x^1$};
  \draw[curva2,desplazada] (0,5.65685) -- (2.82843,0) node[etiqueta,above right,text=tceeazul] {comp.};
  % Curvas de indiferencia
  \draw[curva1] plot[domain=0.78:4.4,samples=80,smooth] (\x,{4/\x}) node[etiqueta,right] {$U^0$};
  \draw[curvanueva] plot[domain=0.4:4.4,samples=80,smooth] (\x,{2/\x}) node[etiqueta,right] {$U^1$};
  % Puntos
  \draw[guia] (\xA,\yA) -- (\xA,-0.15);
  \draw[guia] (\xB,\yB) -- (\xB,-0.15);
  \draw[guia] (\xC,\yC) -- (\xC,-0.15);
  \node[punto] at (\xA,\yA) {}; \node[etiqueta,above right] at (\xA,\yA) {$A$};
  \node[punto] at (\xB,\yB) {}; \node[etiqueta,above right] at (\xB,\yB) {$B$};
  \node[punto] at (\xC,\yC) {}; \node[etiqueta,above left] at (\xC,\yC) {$C$};
  \node[below,font=\footnotesize] at (\xC,-0.1) {$x_2$};
  \node[below,font=\footnotesize] at (\xA,-0.1) {$x_1$};
  % Efectos
  \draw[flecha,tceeazul] (\xA,-0.85) -- (\xB,-0.85) node[midway,below,font=\footnotesize,text=tceeazul] {ES};
  \draw[flecha,tceeverde] (\xB,-0.85) -- (\xC,-0.85) node[midway,below,font=\footnotesize,text=tceeverde] {ER};
  \draw[flecha] (\xA,-1.6) -- (\xC,-1.6) node[midway,below,font=\footnotesize,text=tceerojo] {efecto total};
\end{tikzpicture}
\end{document}
